\section{Background}
\label{background}

A \emph{formal specification} mathematically describes the expected properties of a program. 
Unlike testing, which checks a finite set of inputs, 
formal specifications can quantify over all possible inputs, 
enabling them to provide guarantees over all possible program behaviors.
\emph{Formal verification} is the process of mathematically proving that a program satisfies its specification.
Thus, it shifts the burden from ensuring correctness of an entire implementation 
(which may have complex performance optimizations)
to ensuring correctness of a higher-level mathematical specification, 
an easier task. 

% typically using a verification tool to ensure proof correctness.

One state-of-the-art verification tool is Verus~\cite{Verus}, a semi-automated Rust verifier which 
has already been used at the intersection
of AI and verification \cite{AutoVerus,alphaverus,verussage,RAGVerus,vericode}. 
Verus is written in Rust and uses Rust-style annotations. 
Verus divides into three modes: specification, proof, and executable. 
Specification and proof code is deleted during compilation, so there is no performance cost from adding annotations.
Need to explain proof annotations.
With Verus, programmers can annotate their code with logical pre- and post-conditions, loop invariants, and other proof obligations.
Verus will then check if the code satisfies these requirements. 
% Note that Verus's verification result is only as good as the specification quality: 
% if the specification is incorrect or fails to capture all possible cases, 
% implementations verified against such specifications will also lack full guarantees.
% And since Verus specifications are checkable and compositional, they can serve as unambiguous descriptions of program intent---beyond what unit tests can capture. 
% Verus~\cite{Verus} has already been used in projects at the intersection
% of AI and verification \cite{AutoVerus,alphaverus,verussage,RAGVerus,vericode}. 